\section{VSEPR}

\subsection{Formální náboj}

	\begin{itemize}
		\item Rozdíl mezi počtem valenčních elektronů ve volném atomu a valenčních elektronů ve vázaném atomu.
		\item Záporný náboj je umístěn na nejelektronegativnějším atomu.
		\item Součet formálních nábojů všech atomů v molekule je roven jejímu náboji.
		\item \ce{H3O^+}: H: 0; O: 6-5=+1
		\item \ce{CH3^+}: H: 0; C: 4-3=+1
	\end{itemize}
	\begin{center}
		\adjincludegraphics[width=80mm]{img/formalni_naboj.png}
	\end{center}

\subsection{Rezonanční struktury}
\begin{itemize}
	\item Popisují polohu elektronů v molekulách.
	\item Vyjadřují jednotlivé limitní stavy.
\end{itemize}
\begin{center}
	\schemestart
	\rotatebox{30}{\chemfig{*6(=-=-=-)}}\arrow{<=>}\rotatebox{30}{\chemfig{*6(=-=-=-)}}
	\schemestop
\end{center}
\scalebox{0.7}{
	\schemestart
	\chemfig{S(-[:0]O^{\scriptstyle -})(=[:90]O)(-[:180]^{\scriptstyle -}O)(=[:270]O)}\arrow{<=>}
	\chemfig{S(-[:0]O^{\scriptstyle -})(=[:90]O)(=[:180]O)(-[:270]\chembelow{O}{\scriptstyle -})}\arrow{<=>}
	\chemfig{S(=[:0]O)(-[:90]\chemabove{O}{\scriptstyle -})(=[:180]O)(-[:270]\chembelow{O}{\scriptstyle -})}\arrow{<=>}
	\chemfig{S(=[:0]O)(=[:90]O)(-[:180]^{\scriptstyle -}O)(-[:270]\chembelow{O}{\scriptstyle -})}
	\schemestop
}

\subsection{Molekulové orbitaly}
\begin{itemize}
	\item Teorie LCAO-MO - Linear Combination of Atomic Orbitals - Molecular Orbital.
	\item Molekulové orbitaly vznikají lineární kombinací atomových orbitalů.
	\item Kombinací dvou AO vznikají dva MO - vazebný a protivazebný. Protivazebné orbitaly se označují hvězdičkou, např. $\sigma^*$.
	\item Aby byl překryv úspěšný musí mít vlnové funkce orbitalů v místě překryvu stejná znaménka.
	\item Protivazebný orbital má o jednu nodální plochu více než vazebný a pokud je obsazen elektronovým párem, snižuje řád vazby o jedna. Obsazený vazebný orbital naopak řád vazby zvyšuje.
	\item Vazba $\sigma$ - vzniká osovým překryvem orbitalů.
	\item Vazba $\pi$ - vzniká bočným překryvem orbitalů, je přítomna v násobných vazbách.
	\item Vazba $\delta$ - vzniká překryvem všech čtyř laloků d-orbitalu, je přítomna ve čtverné vazbě např. v \href{http://pubs.acs.org/doi/abs/10.1021/ic50025a015}{\ce{[Re2Cl8]^{2-}}}.
\end{itemize}

\subsection{Řád vazby}
\begin{itemize}
	\item Řád vazby popisuje počet elektronových párů, které tvoří vazbu mezi atomy.
	\item Lze jej odvodit z Lewisovského vzorce molekuly nebo z diagramu MO.
\end{itemize}

\begin{center}
	\chemfig{H-[:52.24]\charge{45=\:,135=\:}{O}-[::-104.48]H}
	\raisebox{3.5ex}{\chemfig{\charge{90=\:,270=\:}{O}=C=\charge{90=\:,270=\:}{O}}}
\end{center}

\begin{itemize}
	\item Řád vazby lze spočítat z počtu elektronů ve vazebných a protivazebných orbitalech.
	\item $RV = \frac{vazebne\ elektrony - protivazebne\ elektrony}{2}$
	
	\item Neobsazené molekulové orbitaly neovlivňují ani řád vazby, ani energii systému.
\end{itemize}

\begin{itemize}
	\item Nepřechodné prvky se snaží vytvářet chemické vazby tak, aby měly ve valenční slupce osm elektronů, čímž dosáhnou na elektronovou konfiguraci vzácného plynu.
	\item Elektrony, které atomy sdílí v kovalentních vazbách se započítávají pro každý atom zvlášť. Např. v molekule \ce{CO2} jsou kyslíky obklopeny čtyřmi nevazebnými elektrony a~čtyři vazebnými, centrální uhlík je pak obklopen celkem osmi elektrony ze čtyř kovalentních vazeb.
\end{itemize}

\begin{center}
	\chemfig{{\charge{90=\:,270=\:}{O}=C=\charge{90=\:,270=\:}{O}}}
\end{center}

\begin{itemize}
	\item Existuje několik výjimek z oktetového pravidla, u nekovů z třetí a~vyšší periody se setkáváme s tzv. \textit{elektronovým dodecetem}, kdy má atom ve valenční slupce 12 elektronů. Toho může dosáhnout díky nezaplněným d-orbitalům.
	\begin{itemize}
		\item Příkladem je molekula \ce{ICl5}, kde má jód dva nevazebné elektrony a~celkem 10 elektronů z pěti kovalentních vazeb.
		\item Podobně síra v molekule \ce{SF6} má ve valenční slupce celkem 12~elektronů ze šesti kovalentních vazeb \ce{S-F}.
	\end{itemize}
\end{itemize}

\subsection{Hybridizace}
\begin{itemize}
	\item Hybridizace atomových orbitalů --- proces energetického mísení a směrového vyrovnání atomových orbitalů daného atomu.
	\item Počet hybridních orbitalů odpovídá počtu mísených atomových orbitalů.
\end{itemize}
\begin{tabular}{|l|l|}
	\hline
	\textbf{Hybridizace} & \textbf{Geometrie molekuly} \\
	\hline
	sp & lineární \\
	\hline
	sp$^2$ & rovnostranný trojúhelník \\
	\hline
	sp$^3$ & tetraedr \\
	\hline
	d$^2$sp$^3$ & oktaedr \\
	\hline
	dsp$^2$ & čtverec \\
	\hline
	\multirow{2}{*}{dsp$^3$} & trigonální bipyramida \\
	& čtvercová pyramida \\
	\hline
\end{tabular}

\subsection{VSEPR}

\begin{itemize}
	\item \textbf{V}alence \textbf{S}hell \textbf{E}lectron \textbf{P}air \textbf{R}epulsion
	\item Tvar molekuly určíme na základě rozmístění elektronových párů v okolí centrálního atomu tak, aby jejich vzájemné odpuzování bylo co nejmenší.
	\item Tento model je vhodný převážně pro sloučeniny nepřechodných prvků.
	\item Uvažujeme pouze nevazebné elektronové páry - $n$ a vazebné elektronové páry $\sigma$.
	\item \textbf{Základní pravidla VSEPRu}
	\begin{enumerate}
		\item Elektronové páry centrálního atomu se v prostoru rozmístí tak, aby byly co nejdále od sebe a měly minimální energii.
		\item Nevazebný elektronový pár odpuzuje ostatní elektronové páry nejvíce, odpuzování  vazebných elektronových párů je slabší a klesá v pořadí:
		\begin{itemize}
			\item trojná vazba $>$ dvojná vazba $>$ jednoduchá vazba.
		\end{itemize}
		\item Tvar molekuly je dán pouze polohou vazebných elektronových párů.
	\end{enumerate}
\end{itemize}

\begin{tabular}{|l|l|l|l|l|}
	\hline
	\textbf{Počet el. párů} & \textbf{Notace} & \textbf{Tvar} & \textbf{Úhel} & \textbf{Příklad} \\\hline
	2 & AX2 & Lineární & 180 & \ce{CO2} \\\hline
	\multirow{3}{*}{3} & AX3 & Rovnostranný trojúhelník & 120 & \ce{BCl3} \\\cline{2-5}
	& AX2E & Lomenný & <120 & \ce{SO2} \\\cline{2-5}
	& AXE2 & Lineární & 180 & \ce{O2} \\\hline
	\multirow{3}{*}{4} & AX4 & Tetraedr & 109,5 & \ce{CH4}, \ce{SO$^{2-}_4$} \\\cline{2-5}
	& AX3E & Trigonální pyramida & <109,5 & \ce{NH3} \\\cline{2-5}
	& AX2E2 & Lomenný & <109,5 & \ce{H2O} \\\hline
	\multirow{4}{*}{5} & AX5 & Trigonální bipyramida & 90 a 120 & \ce{AsF5}, \ce{PCl5} \\\cline{2-5}
	& AX4E & Houpačka & <90 a <120 & \ce{SeH4} \\\cline{2-5}
	& AX3E2 & Tvar T & 90 & \ce{ICl3} \\\cline{2-5}
	& AX2E3 & Lineární & 180 & \ce{IF$^-_2$} \\\hline
	\multirow{3}{*}{6} & AX6 & Oktaedr & 90 & \ce{SF6}, \ce{SCl6} \\\cline{2-5}
	& AX5E & Čtvercová pyramida & 90 & \ce{IF5} \\\cline{2-5}
	& AX4E2 & Čtverec & 90 & \ce{XeF4} \\\hline
\end{tabular}

\clearpage

\subsection{Symetrie}
\begin{itemize}
	\item \textbf{Operace symetrie} - geometrická operace, jejímž provedením dostaneme objekt do polohy nerozlišitelné od výchozí.
	\item \textbf{Prvek symetrie} - body, jejichž poloha se v průběhu provádění operace symetrie nemění.
	\item U molekul existuje pět prvků symetrie.
\end{itemize}

\begin{tabular}{|l|l|l|}
	\hline
	\textbf{Operace symetrie} & \textbf{Symbol} & \textbf{Prvek symetrie} \\
	\hline
	Identita & E & Celý objekt \\
	\hline
	Rotace & $C_n$ & Rotační osa \\
	\hline
	Zrcadlení & $\sigma$ & Rovina symetrie \\
	\hline
	Inverze & i & Střed symetrie \\
	\hline
	Nevlastní osa & $S_n$ & Rotačně-reflexní osa \\
	\hline
\end{tabular}

\subsection{Dipólový moment}
\begin{itemize}
	\item Vektor popisující rozložení elektrického náboje v molekule.
	\item Výsledný dipólmoment získáme vektorovým součtem dipólmomentů jednotlivých vazeb.
\end{itemize}
\begin{center}
	\begin{picture}(0,0)
		\color{red}
		\put(0,0){O}
		\color{blue}
		\put(23,-30){H}
		\put(-23,-30){H}
		\color{black}
		\thicklines
		\put(4,-1){\vector(-1,-1){20}}
		\put(4,-1){\vector(1,-1){20}}
		\put(4,-1){\vector(0,-1){40}}
		\thinlines
		\put(24,-21){\line(-1,-1){20}}
		\put(-16,-21){\line(1,-1){20}}
	\end{picture}
\end{center}